\paragraph{\trFunctoriality。}\label{sub:twist-functoriality}
设 $\scr{F}\colon\scr{A}\to\scr{B}$ 是\trACategory{}之间的一个%
\trCohomology\trFullyFaithful{}的\trFunctor。取 $Y_0$、$Y_1\in\Ob\scr{A}$。

\begin{lemma}\label{thm:twist-functorial-isomorphism}
设扭转 $T_{Y_0}Y_1$ 与 $T_{\scr{F}(Y_0)}\scr{F}(Y_1)$ 二者都存在。那么在
$H^0(\scr{B})$ 中有一个\trCanonical{}的\trIsomorphism
\begin{equation}\label{eq:functorial-twist-isomorphism}
\scr{F}(T_{Y_0}Y_1)\cong T_{\scr{F}(Y_0)}\scr{F}(Y_1).
\end{equation}
\end{lemma}

\begin{proof}
我们先从论证的一个较简单的初步版本入手。据\cref{thm:triangle-completion}，%
\trExactTriangle{} \eqref{eq:twisted-yoneda-triangle} 在 $H^0(\scr{A})$ 中把
$T_{Y_0}(Y_1)$ 刻画到相差\trIsomorphism{}的程度。由这一点与%
\cref{thm:afunctor-maps-exact-triangles}容易推出，%
\eqref{eq:functorial-twist-isomorphism} 中的两个\trObject{}是\trIsomorphic{}的。
然而，这并不能给出\trIsomorphism{}的一个\trCanonical{}的选取；要得到更精细的
陈述，还得稍微细心一些来论证。

沿着与\eqref{eq:yoneda-pullback-transformation} 相同的思路，对任意的
$Y\in\Ob\scr{A}$ 与 $\scr{M}\in\Ob(\mathit{mod}(\scr{B}))$，都有一个%
\trCanonical{}的\trModule{}\trHomomorphism
\begin{equation*}
s\colon\scr{T}_Y(\scr{F}^*\scr{M})\longrightarrow
\scr{F}^*\bigl(\scr{T}_{\scr{F}(Y)}\scr{M}\bigr).
\end{equation*}
具体地说，$s^1$ 把 $(c_0\otimes b,c_1)$ 映为
$\bigl((-1)^{|c_0|+|b|-1}c_0\otimes\scr{F}^1(b),\,(-1)^{|c_1|}c_1\bigr)$；
而当 $d\ge 2$ 时，$s^d$ 把
$\bigl((c_0\otimes b,c_1),a_{d-1},\ldots,a_1\bigr)$ 映为
$\bigl((-1)^{\Delta}c_0\otimes\scr{F}^d(b,a_{d-1},\ldots,a_1),\,0\bigr)$，
其中 $\Delta=|a_1|+\cdots+|a_{d-1}|+|b|+|c_0|-d$。看 $d=1$ 的情形便知，$[s]$
是 $H^0(\mathit{mod}(\scr{A}))$ 中的一个\trIsomorphism。把它应用于
$Y=Y_0$、$\scr{M}=(\scr{J}_{\scr{B}}\circ\scr{F})(Y_1)$ 便得到
\begin{equation}\label{eq:twist-pullback-yoneda}
\scr{T}_{Y_0}(\scr{F}^*\circ\scr{J}_{\scr{B}}\circ\scr{F})(Y_1)\cong
\scr{F}^*\bigl(\scr{T}_{\scr{F}(Y_0)}(\scr{J}_{\scr{B}}\circ\scr{F})(Y_1)\bigr).
\end{equation}
据\eqref{eq:yoneda-naturality-diagram}，\eqref{eq:twist-pullback-yoneda}
的左端\trCanonical{}地\trIsomorphic{}于 $\scr{T}_{Y_0}(\scr{J}_{\scr{A}}(Y_1))$；
而由 $T_{Y_0}$ 的
定义，后者又\trCanonical{}地\trIsomorphic{}于 $\scr{J}_{\scr{A}}(T_{Y_0}(Y_1))$，
从而\trIsomorphic{}于 $(\scr{F}^*\circ\scr{J}_{\scr{B}}\circ\scr{F})(T_{Y_0}Y_1)$。
把 $T_{\scr{F}(Y_0)}$ 的定义应用于右端即可发现，它\trCanonical{}地%
\trIsomorphic{}于 $(\scr{F}^*\circ\scr{J}_{\scr{B}})
(T_{\scr{F}(Y_0)}\scr{F}(Y_1))$。结论是：\eqref{eq:functorial-twist-isomorphism}
中的两个\trObject{}在经过 $\scr{F}^*\circ\scr{J}_{\scr{B}}$ 之后变得%
\trCanonical{}地\trIsomorphic{}了。再利用\eqref{eq:yoneda-naturality-diagram}
可知，$\scr{F}^*\circ\scr{J}_{\scr{B}}$ 在 $\scr{F}$ 的\trImage{}上是%
\trCohomology\trFullyFaithful{}的，因而在 $H^0(\scr{B})$ 中由那些与某个
$\scr{F}(X)$ \trIsomorphic{}的\trObject{}组成的稍大一类上也是如此。既然
\eqref{eq:functorial-twist-isomorphism} 的两端都属于这一类，我们便得到二者之间
一个\trCanonical{}的\trIsomorphism。
\end{proof}

这样构造出的\trIsomorphism{}关于 $Y_1$ 是自然的（事实上，只要再多费些力气，
便可以把它扩张成\trFunctor\trCategory{} $H^0(\mathit{fun}(\scr{A},\scr{B}))$ 中
$\scr{F}\circ T_{Y_0}$ 与 $T_{\scr{F}(Y_0)}\circ\scr{F}$ 之间的一个%
\trCanonical{}的\trIsomorphism）。它还与\trCanonical{}映射
$Y_1\to T_{Y_0}Y_1$ 及 $\scr{F}(Y_1)\to T_{\scr{F}(Y_0)}\scr{F}(Y_1)$
\trCompatible。如果进一步假设我们的 \trACategory{}对有限\trDirectSum{}与%
\trSuspension{}封闭，并且相关的 $\Hom$ 空间是有限维的，那么\trExactTriangle
\eqref{eq:twisted-yoneda-triangle} 在 $H(\scr{F})$ 下的\trImage{}就与
$\scr{F}(Y_0)$、$\scr{F}(Y_1)$ 这两个\trObject{}的相应三角形\trCanonical{}地%
\trIsomorphic{}了。
